
\documentclass{article}
\usepackage{colortbl}
\usepackage{makecell}
\usepackage{multirow}
\usepackage{supertabular}

\begin{document}

\newcounter{utterance}

\centering \large Interaction Transcript for game `clean\_up', experiment `0\_easy\_5obj\_cs', episode 4 with mistral{-}large{-}3{-}azure.
\vspace{24pt}

{ \footnotesize  \setcounter{utterance}{1}
\setlength{\tabcolsep}{0pt}
\begin{supertabular}{c@{$\;$}|p{.15\linewidth}@{}p{.15\linewidth}p{.15\linewidth}p{.15\linewidth}p{.15\linewidth}p{.15\linewidth}}
    \# & $\;$A & \multicolumn{4}{c}{Game Master} & $\;\:$B\\
    \hline

    \theutterance \stepcounter{utterance}  
    & & \multicolumn{4}{p{0.6\linewidth}}{
        \cellcolor[rgb]{0.9,0.9,0.9}{
            \makecell[{{p{\linewidth}}}]{
                \texttt{\tiny{[P1$\langle$GM]}}
                \texttt{Jsem váš herní mistr a hrajete společnou hru s následující mřížkou jako herním plánem:} \\
\\ 
\texttt{```} \\
\texttt{ 1234567} \\
\texttt{╔═╤═════╗} \\
\texttt{║◌│◌◌◌◌◌║ 1} \\
\texttt{║◌│I◌◌◌◌║ 2} \\
\texttt{║C│◌◌◌◌◌║ 3} \\
\texttt{║H│◌W◌◌◌║ 4} \\
\texttt{╟─┼───┐◌║ 5} \\
\texttt{║◌│◌◌◌├─╢ 6} \\
\texttt{║◌│◌◌◌│T║ 7} \\
\texttt{╚═╧═══╧═╝} \\
\\ 
\texttt{```} \\
\\ 
\texttt{* Horní okraj zobrazuje souřadnice x rostoucí doprava a pravý okraj souřadnice y rostoucí dolů.} \\
\texttt{* Na vaší mřížce jsou náhodně umístěny následující objekty: 'W', 'I', 'T', 'C', 'H'.} \\
\\ 
\\ 
\texttt{Druhý hráč vidí variantu herní desky, kde jsou objekty umístěny na různých náhodných místech. Vy nevidíte desku druhého hráče a on nevidí tu vaši.} \\
\\ 
\\ 
\texttt{**Cíl:**} \\
\\ 
\texttt{Oba hráči musí přesunout objekty na svém pozadí tak, aby stejné objekty skončily na stejných souřadnicích. Musíte komunikovat s druhým hráčem, abyste se dohodli na společné cílové konfiguraci.} \\
\\ 
\\ 
\texttt{**Pravidla:**} \\
\\ 
\texttt{* V každém tahu můžete poslat přesně jeden z následujících dvou příkazů:} \\
\texttt{ 1. `ŘEKNI: <ZPRÁVA>`: odeslání zprávy (vše až do dalšího zalomení řádku) druhému hráči. Přepošlu ji vašemu partnerovi.} \\
\texttt{2. `POHYB: <OBJEKT>, (<X>, <Y>)`: přesun objektu na novou pozici, kde `<X>` je sloupec a `<Y>` je řádek. Budu vás informovat, zda byl váš přesun úspěšný, nebo ne.} \\
\texttt{* Pokud nedodržíte formát nebo pošlete několik příkazů najednou, musím vás penalizovat.} \\
\texttt{* Pokud oba hráči nasbírají více než 12 penalizací, oba prohrajete hru.} \\
\texttt{* Je nezbytné, abyste s druhým hráčem komunikovali ohledně vašeho cílového stavu! *Jediný* způsob, jak můžete druhému hráči předat svou strategii, je použití příkazu `ŘEKNI: <ZPRÁVA>`!} \\
\\ 
\\ 
\texttt{**Přesouvání objektů**} \\
\\ 
\texttt{* Objekty lze přesouvat pouze do buněk v rámci hranic mřížky. Cílová buňka musí být prázdná, tj. musí obsahovat pouze symbol '◌'.} \\
\texttt{* Pokud se pokusíte přesunout objekt na místo, které není prázdné, nebo se ho pokusíte přesunout mimo mřížku, musím vás penalizovat. Dostanete další pokus.} \\
\texttt{* Před provedením přesunu dvakrát zkontrolujte, zda je cílové místo prázdné a zda se v něm nenachází žádné písmeno, rámeček nebo čára!} \\
\\ 
\\ 
\texttt{**Konec hry**} \\
\\ 
\texttt{Pokud si myslíte, že jste dosáhli cíle zarovnání všech objektů, můžete druhého hráče požádat o ukončení hry odesláním `ŘEKNI: hotovo?`. Pokud vás druhý hráč požádá o ukončení hry a vy odpovíte `ŘEKNI: hotovo!`, hra skončí.} \\
\\ 
\texttt{Oba hráči vyhrají, pokud hra skončí během 20 kol, přičemž jedno kolo je definováno jako odeslání platného příkazu oběma hráči.} \\
\\ 
\texttt{**Bodování:**} \\
\\ 
\texttt{Čím blíže jsou stejné objekty na obou herních deskách, tím více bodů získáte. Penalizace vám body snižují. Podaří se vám překonat rekord?} \\
\\ 
\\ 
\texttt{Pošlete prosím zprávu druhému hráči a začněte hru!} \\
            }
        }
    }
    & \\ \\

    \theutterance \stepcounter{utterance}  
    & \multicolumn{4}{p{0.6\linewidth}}{
        \cellcolor[rgb]{0.9,0.9,0.9}{
            \makecell[{{p{\linewidth}}}]{
                \texttt{\tiny{[P1$\rangle$GM]}}
                \texttt{Začneme hru komunikací s druhým hráčem, abychom se dohodli na společné cílové konfiguraci. Navrhnu strategii, jak objekty uspořádat, a budu se snažit minimalizovat počet tahů.} \\
\\ 
\texttt{**První zpráva pro druhého hráče:**} \\
\texttt{```} \\
\texttt{ŘEKNI: Ahoj! Navrhuji uspořádat objekty do pravidelného vzoru pro snadné zarovnání. Například:} \\
\texttt{{-} W na (3,4) → (2,2)} \\
\texttt{{-} I na (3,2) → (3,2) (zůstane na místě)} \\
\texttt{{-} T na (6,7) → (4,4)} \\
\texttt{{-} C na (2,3) → (5,5)} \\
\texttt{{-} H na (2,4) → (6,6)} \\
\texttt{Souhlasíš? Pokud máš jiný návrh, dej vědět!} \\
\texttt{```} \\
\\ 
\texttt{Tento návrh rozmisťuje objekty do diagonály nebo jiného snadno zapamatovatelného vzoru. Počkáme na odpověď druhého hráče a podle toho upravíme strategii. Pokud souhlasí, můžeme začít přesuny. Pokud ne, budeme diskutovat dál.} \\
\\ 
\texttt{**Poznámka:** Pokud druhý hráč navrhne jinou konfiguraci, budeme se muset přizpůsobit. Důležité je najít společný cíl co nejrychleji.} \\
            }
        }
    }
    & & \\ \\

    \theutterance \stepcounter{utterance}  
    & & & \multicolumn{2}{p{0.3\linewidth}}{
        \cellcolor[rgb]{0.9,0.9,0.9}{
            \makecell[{{p{\linewidth}}}]{
                \texttt{\tiny{[GM$|$GM]}}
                \texttt{Začneme hru komunikací s druhým hráčem, abychom se dohodli na společné cílové konfiguraci. Navrhnu strategii, jak objekty uspořádat, a budu se snažit minimalizovat počet tahů.} \\
\\ 
\texttt{**První zpráva pro druhého hráče:**} \\
\texttt{```} \\
\texttt{ŘEKNI: Ahoj! Navrhuji uspořádat objekty do pravidelného vzoru pro snadné zarovnání. Například:} \\
\texttt{{-} W na (3,4) → (2,2)} \\
\texttt{{-} I na (3,2) → (3,2) (zůstane na místě)} \\
\texttt{{-} T na (6,7) → (4,4)} \\
\texttt{{-} C na (2,3) → (5,5)} \\
\texttt{{-} H na (2,4) → (6,6)} \\
\texttt{Souhlasíš? Pokud máš jiný návrh, dej vědět!} \\
\texttt{```} \\
\\ 
\texttt{Tento návrh rozmisťuje objekty do diagonály nebo jiného snadno zapamatovatelného vzoru. Počkáme na odpověď druhého hráče a podle toho upravíme strategii. Pokud souhlasí, můžeme začít přesuny. Pokud ne, budeme diskutovat dál.} \\
\\ 
\texttt{**Poznámka:** Pokud druhý hráč navrhne jinou konfiguraci, budeme se muset přizpůsobit. Důležité je najít společný cíl co nejrychleji.} \\
            }
        }
    }
    & & \\ \\

    \theutterance \stepcounter{utterance}  
    & & & \multicolumn{2}{p{0.3\linewidth}}{
        \cellcolor[rgb]{0.9,0.9,0.9}{
            \makecell[{{p{\linewidth}}}]{
                \texttt{\tiny{[GM$|$GM]}}
                \texttt{Invalid format: head and tail are not empty} \\
\texttt{head: 'Začneme hru komunikací s druhým hráčem, abychom se dohodli na společné cílové konfiguraci. Navrhnu strategii, jak objekty uspořádat, a budu se snažit minimalizovat počet tahů.} \\
\\ 
\texttt{**První zpráva pro druhého hráče:**} \\
\\ 
\texttt{'} \\
\texttt{tail: '} \\
\texttt{{-} W na (3,4) → (2,2)} \\
\texttt{{-} I na (3,2) → (3,2) (zůstane na místě)} \\
\texttt{{-} T na (6,7) → (4,4)} \\
\texttt{{-} C na (2,3) → (5,5)} \\
\texttt{{-} H na (2,4) → (6,6)} \\
\texttt{Souhlasíš? Pokud máš jiný návrh, dej vědět!} \\
\\ 
\\ 
\texttt{Tento návrh rozmisťuje objekty do diagonály nebo jiného snadno zapamatovatelného vzoru. Počkáme na odpověď druhého hráče a podle toho upravíme strategii. Pokud souhlasí, můžeme začít přesuny. Pokud ne, budeme diskutovat dál.} \\
\\ 
\texttt{**Poznámka:** Pokud druhý hráč navrhne jinou konfiguraci, budeme se muset přizpůsobit. Důležité je najít společný cíl co nejrychleji.'} \\
            }
        }
    }
    & & \\ \\

    \theutterance \stepcounter{utterance}  
    & & \multicolumn{4}{p{0.6\linewidth}}{
        \cellcolor[rgb]{0.9,0.9,0.9}{
            \makecell[{{p{\linewidth}}}]{
                \texttt{\tiny{[P1$\langle$GM]}}
                \texttt{Trest: Vaše zpráva nesmí obsahovat nic před příkazem ani po něm!} \\
\\ 
\texttt{Ujistěte se, že vaše odpověď obsahuje pouze `ŘEKNI: <ZPRÁVA>` nebo `POHYB: <OBJEKT>, (<X>, <Y>)` a nic jiného!} \\
\\ 
\texttt{Dohromady jste nasbírali 1 z 12 pokut.} \\
\texttt{Zkuste to prosím znovu!} \\
            }
        }
    }
    & \\ \\

    \theutterance \stepcounter{utterance}  
    & \multicolumn{4}{p{0.6\linewidth}}{
        \cellcolor[rgb]{0.9,0.9,0.9}{
            \makecell[{{p{\linewidth}}}]{
                \texttt{\tiny{[P1$\rangle$GM]}}
                \texttt{```} \\
\texttt{ŘEKNI: Navrhuji umístit W na (2,2), I na (3,2), T na (4,4), C na (5,5), H na (6,6). Souhlasíš?} \\
\texttt{```} \\
            }
        }
    }
    & & \\ \\

    \theutterance \stepcounter{utterance}  
    & & & \multicolumn{2}{p{0.3\linewidth}}{
        \cellcolor[rgb]{0.9,0.9,0.9}{
            \makecell[{{p{\linewidth}}}]{
                \texttt{\tiny{[GM$|$GM]}}
                \texttt{```} \\
\texttt{ŘEKNI: Navrhuji umístit W na (2,2), I na (3,2), T na (4,4), C na (5,5), H na (6,6). Souhlasíš?} \\
\texttt{```} \\
            }
        }
    }
    & & \\ \\

    \theutterance \stepcounter{utterance}  
    & & \multicolumn{4}{p{0.6\linewidth}}{
        \cellcolor[rgb]{0.9,0.9,0.9}{
            \makecell[{{p{\linewidth}}}]{
                \texttt{\tiny{[GM$\rangle$P2]}}
                \texttt{Jsem váš herní mistr a hrajete společnou hru s následující mřížkou jako herním plánem:} \\
\\ 
\texttt{```} \\
\texttt{ 1234567} \\
\texttt{╔═╤═════╗} \\
\texttt{║◌│◌◌◌◌◌║ 1} \\
\texttt{║◌│◌◌◌H◌║ 2} \\
\texttt{║I│◌◌C◌◌║ 3} \\
\texttt{║◌│◌◌T◌◌║ 4} \\
\texttt{╟─┼───┐◌║ 5} \\
\texttt{║◌│◌◌◌├─╢ 6} \\
\texttt{║W│◌◌◌│◌║ 7} \\
\texttt{╚═╧═══╧═╝} \\
\\ 
\texttt{```} \\
\\ 
\texttt{* Horní okraj zobrazuje souřadnice x rostoucí doprava a pravý okraj souřadnice y rostoucí dolů.} \\
\texttt{* Na vaší mřížce jsou náhodně umístěny následující objekty: 'W', 'I', 'T', 'C', 'H'.} \\
\\ 
\\ 
\texttt{Druhý hráč vidí variantu herní desky, kde jsou objekty umístěny na různých náhodných místech. Vy nevidíte desku druhého hráče a on nevidí tu vaši.} \\
\\ 
\\ 
\texttt{**Cíl:**} \\
\\ 
\texttt{Oba hráči musí přesunout objekty na svém pozadí tak, aby stejné objekty skončily na stejných souřadnicích. Musíte komunikovat s druhým hráčem, abyste se dohodli na společné cílové konfiguraci.} \\
\\ 
\\ 
\texttt{**Pravidla:**} \\
\\ 
\texttt{* V každém tahu můžete poslat přesně jeden z následujících dvou příkazů:} \\
\texttt{ 1. `ŘEKNI: <ZPRÁVA>`: odeslání zprávy (vše až do dalšího zalomení řádku) druhému hráči. Přepošlu ji vašemu partnerovi.} \\
\texttt{2. `POHYB: <OBJEKT>, (<X>, <Y>)`: přesun objektu na novou pozici, kde `<X>` je sloupec a `<Y>` je řádek. Budu vás informovat, zda byl váš přesun úspěšný, nebo ne.} \\
\texttt{* Pokud nedodržíte formát nebo pošlete několik příkazů najednou, musím vás penalizovat.} \\
\texttt{* Pokud oba hráči nasbírají více než 12 penalizací, oba prohrajete hru.} \\
\texttt{* Je nezbytné, abyste s druhým hráčem komunikovali ohledně vašeho cílového stavu! *Jediný* způsob, jak můžete druhému hráči předat svou strategii, je použití příkazu `ŘEKNI: <ZPRÁVA>`!} \\
\\ 
\\ 
\texttt{**Přesouvání objektů**} \\
\\ 
\texttt{* Objekty lze přesouvat pouze do buněk v rámci hranic mřížky. Cílová buňka musí být prázdná, tj. musí obsahovat pouze symbol '◌'.} \\
\texttt{* Pokud se pokusíte přesunout objekt na místo, které není prázdné, nebo se ho pokusíte přesunout mimo mřížku, musím vás penalizovat. Dostanete další pokus.} \\
\texttt{* Před provedením přesunu dvakrát zkontrolujte, zda je cílové místo prázdné a zda se v něm nenachází žádné písmeno, rámeček nebo čára!} \\
\\ 
\\ 
\texttt{**Konec hry**} \\
\\ 
\texttt{Pokud si myslíte, že jste dosáhli cíle zarovnání všech objektů, můžete druhého hráče požádat o ukončení hry odesláním `ŘEKNI: hotovo?`. Pokud vás druhý hráč požádá o ukončení hry a vy odpovíte `ŘEKNI: hotovo!`, hra skončí.} \\
\\ 
\texttt{Oba hráči vyhrají, pokud hra skončí během 20 kol, přičemž jedno kolo je definováno jako odeslání platného příkazu oběma hráči.} \\
\\ 
\texttt{**Bodování:**} \\
\\ 
\texttt{Čím blíže jsou stejné objekty na obou herních deskách, tím více bodů získáte. Penalizace vám body snižují. Podaří se vám překonat rekord?} \\
\\ 
\texttt{Druhý hráč zahájil hru odesláním této zprávy:} \\
\\ 
\texttt{"Navrhuji umístit W na (2,2), I na (3,2), T na (4,4), C na (5,5), H na (6,6). Souhlasíš?"} \\
\\ 
\texttt{Jaký je váš první příkaz?} \\
            }
        }
    }
    & \\ \\

    \theutterance \stepcounter{utterance}  
    & & & \multicolumn{4}{p{0.6\linewidth}}{
        \cellcolor[rgb]{0.9,0.9,0.9}{
            \makecell[{{p{\linewidth}}}]{
                \texttt{\tiny{[GM$\langle$P2]}}
                \texttt{Děkuji za vysvětlení a za zahájení komunikace druhým hráčem! Navržená cílová konfigurace vypadá rozumně, ale nejprve bych chtěl ověřit aktuální pozice mých objektů a zkontrolovat, zda jsou některé z nich již na navržených cílových pozicích nebo blízko nich.} \\
\\ 
\texttt{Zde je můj první příkaz – požádám o potvrzení navržené konfigurace a zároveň se zeptám na aktuální pozice objektů druhého hráče, abychom mohli lépe koordinovat přesuny:} \\
\\ 
\texttt{```} \\
\texttt{ŘEKNI: Souhlasím s navrženou konfigurací jako cílovou. Můžeš mi prosím sdělit aktuální pozice tvých objektů? Moje jsou: W na (1,7), I na (1,3), T na (4,4), C na (5,3), H na (6,2). Díky!} \\
\texttt{```} \\
\\ 
\texttt{Tento příkaz:} \\
\texttt{1. Potvrzuje souhlas s cílovou konfigurací.} \\
\texttt{2. Poskytuje mé aktuální pozice objektů pro lepší koordinaci.} \\
\texttt{3. Ptá se na pozice objektů druhého hráče, abychom mohli efektivně plánovat přesuny.} \\
            }
        }
    }
     \\ \\

    \theutterance \stepcounter{utterance}  
    & & & \multicolumn{2}{p{0.3\linewidth}}{
        \cellcolor[rgb]{0.9,0.9,0.9}{
            \makecell[{{p{\linewidth}}}]{
                \texttt{\tiny{[GM$|$GM]}}
                \texttt{Děkuji za vysvětlení a za zahájení komunikace druhým hráčem! Navržená cílová konfigurace vypadá rozumně, ale nejprve bych chtěl ověřit aktuální pozice mých objektů a zkontrolovat, zda jsou některé z nich již na navržených cílových pozicích nebo blízko nich.} \\
\\ 
\texttt{Zde je můj první příkaz – požádám o potvrzení navržené konfigurace a zároveň se zeptám na aktuální pozice objektů druhého hráče, abychom mohli lépe koordinovat přesuny:} \\
\\ 
\texttt{```} \\
\texttt{ŘEKNI: Souhlasím s navrženou konfigurací jako cílovou. Můžeš mi prosím sdělit aktuální pozice tvých objektů? Moje jsou: W na (1,7), I na (1,3), T na (4,4), C na (5,3), H na (6,2). Díky!} \\
\texttt{```} \\
\\ 
\texttt{Tento příkaz:} \\
\texttt{1. Potvrzuje souhlas s cílovou konfigurací.} \\
\texttt{2. Poskytuje mé aktuální pozice objektů pro lepší koordinaci.} \\
\texttt{3. Ptá se na pozice objektů druhého hráče, abychom mohli efektivně plánovat přesuny.} \\
            }
        }
    }
    & & \\ \\

    \theutterance \stepcounter{utterance}  
    & & & \multicolumn{2}{p{0.3\linewidth}}{
        \cellcolor[rgb]{0.9,0.9,0.9}{
            \makecell[{{p{\linewidth}}}]{
                \texttt{\tiny{[GM$|$GM]}}
                \texttt{Invalid format: head and tail are not empty} \\
\texttt{head: 'Děkuji za vysvětlení a za zahájení komunikace druhým hráčem! Navržená cílová konfigurace vypadá rozumně, ale nejprve bych chtěl ověřit aktuální pozice mých objektů a zkontrolovat, zda jsou některé z nich již na navržených cílových pozicích nebo blízko nich.} \\
\\ 
\texttt{Zde je můj první příkaz – požádám o potvrzení navržené konfigurace a zároveň se zeptám na aktuální pozice objektů druhého hráče, abychom mohli lépe koordinovat přesuny:} \\
\\ 
\\ 
\texttt{'} \\
\texttt{tail: '} \\
\\ 
\\ 
\texttt{Tento příkaz:} \\
\texttt{1. Potvrzuje souhlas s cílovou konfigurací.} \\
\texttt{2. Poskytuje mé aktuální pozice objektů pro lepší koordinaci.} \\
\texttt{3. Ptá se na pozice objektů druhého hráče, abychom mohli efektivně plánovat přesuny.'} \\
            }
        }
    }
    & & \\ \\

    \theutterance \stepcounter{utterance}  
    & & \multicolumn{4}{p{0.6\linewidth}}{
        \cellcolor[rgb]{0.9,0.9,0.9}{
            \makecell[{{p{\linewidth}}}]{
                \texttt{\tiny{[GM$\rangle$P2]}}
                \texttt{Trest: Vaše zpráva nesmí obsahovat nic před příkazem ani po něm!} \\
\\ 
\texttt{Ujistěte se, že vaše odpověď obsahuje pouze `ŘEKNI: <ZPRÁVA>` nebo `POHYB: <OBJEKT>, (<X>, <Y>)` a nic jiného!} \\
\\ 
\texttt{Dohromady jste nasbírali 2 z 12 pokut.} \\
\texttt{Zkuste to prosím znovu!} \\
            }
        }
    }
    & \\ \\

    \theutterance \stepcounter{utterance}  
    & & & \multicolumn{4}{p{0.6\linewidth}}{
        \cellcolor[rgb]{0.9,0.9,0.9}{
            \makecell[{{p{\linewidth}}}]{
                \texttt{\tiny{[GM$\langle$P2]}}
                \texttt{Rozumím, budu dodržovat formát. Zde je opravený příkaz:} \\
\\ 
\texttt{```} \\
\texttt{ŘEKNI: Souhlasím s navrženou konfigurací. Jaké jsou aktuální pozice tvých objektů?} \\
\texttt{```} \\
            }
        }
    }
     \\ \\

    \theutterance \stepcounter{utterance}  
    & & & \multicolumn{2}{p{0.3\linewidth}}{
        \cellcolor[rgb]{0.9,0.9,0.9}{
            \makecell[{{p{\linewidth}}}]{
                \texttt{\tiny{[GM$|$GM]}}
                \texttt{Rozumím, budu dodržovat formát. Zde je opravený příkaz:} \\
\\ 
\texttt{```} \\
\texttt{ŘEKNI: Souhlasím s navrženou konfigurací. Jaké jsou aktuální pozice tvých objektů?} \\
\texttt{```} \\
            }
        }
    }
    & & \\ \\

    \theutterance \stepcounter{utterance}  
    & & & \multicolumn{2}{p{0.3\linewidth}}{
        \cellcolor[rgb]{0.9,0.9,0.9}{
            \makecell[{{p{\linewidth}}}]{
                \texttt{\tiny{[GM$|$GM]}}
                \texttt{Invalid format: head is not empty: 'Rozumím, budu dodržovat formát. Zde je opravený příkaz:} \\
\\ 
\\ 
\texttt{'} \\
            }
        }
    }
    & & \\ \\

    \theutterance \stepcounter{utterance}  
    & & \multicolumn{4}{p{0.6\linewidth}}{
        \cellcolor[rgb]{0.9,0.9,0.9}{
            \makecell[{{p{\linewidth}}}]{
                \texttt{\tiny{[GM$\rangle$P2]}}
                \texttt{Trest: Vaše zpráva nesmí obsahovat nic před příkazem!} \\
\\ 
\texttt{Ujistěte se, že vaše odpověď obsahuje pouze `ŘEKNI: <ZPRÁVA>` nebo `POHYB: <OBJEKT>, (<X>, <Y>)` a nic jiného!} \\
\\ 
\texttt{Dohromady jste nasbírali 3 z 12 pokut.} \\
\texttt{Zkuste to prosím znovu!} \\
            }
        }
    }
    & \\ \\

    \theutterance \stepcounter{utterance}  
    & & & \multicolumn{4}{p{0.6\linewidth}}{
        \cellcolor[rgb]{0.9,0.9,0.9}{
            \makecell[{{p{\linewidth}}}]{
                \texttt{\tiny{[GM$\langle$P2]}}
                \texttt{```} \\
\texttt{ŘEKNI: Souhlasím. Jaké máš aktuální pozice?} \\
\texttt{```} \\
            }
        }
    }
     \\ \\

    \theutterance \stepcounter{utterance}  
    & & & \multicolumn{2}{p{0.3\linewidth}}{
        \cellcolor[rgb]{0.9,0.9,0.9}{
            \makecell[{{p{\linewidth}}}]{
                \texttt{\tiny{[GM$|$GM]}}
                \texttt{```} \\
\texttt{ŘEKNI: Souhlasím. Jaké máš aktuální pozice?} \\
\texttt{```} \\
            }
        }
    }
    & & \\ \\

    \theutterance \stepcounter{utterance}  
    & & \multicolumn{4}{p{0.6\linewidth}}{
        \cellcolor[rgb]{0.9,0.9,0.9}{
            \makecell[{{p{\linewidth}}}]{
                \texttt{\tiny{[P1$\langle$GM]}}
                \texttt{Vaše zpráva byla předána druhému hráči.} \\
\\ 
\\ 
\texttt{Právě hrajete kolo 1 z maximálně 20.} \\
\texttt{Dohromady jste nasbírali 3 z 12 pokut.} \\
\texttt{Druhý hráč poslal tuto zprávu:} \\
\\ 
\\ 
\texttt{"Souhlasím. Jaké máš aktuální pozice?"} \\
\\ 
\\ 
\texttt{Jaký je váš další příkaz?} \\
            }
        }
    }
    & \\ \\

    \theutterance \stepcounter{utterance}  
    & \multicolumn{4}{p{0.6\linewidth}}{
        \cellcolor[rgb]{0.9,0.9,0.9}{
            \makecell[{{p{\linewidth}}}]{
                \texttt{\tiny{[P1$\rangle$GM]}}
                \texttt{```} \\
\texttt{ŘEKNI: Moje pozice: W(3,4), I(3,2), T(6,7), C(2,3), H(2,4). Začnu přesouvat W na (2,2).} \\
\texttt{```} \\
\\ 
\texttt{*(Počkáme na potvrzení druhého hráče a pak provedeme přesun W.)*} \\
            }
        }
    }
    & & \\ \\

    \theutterance \stepcounter{utterance}  
    & & & \multicolumn{2}{p{0.3\linewidth}}{
        \cellcolor[rgb]{0.9,0.9,0.9}{
            \makecell[{{p{\linewidth}}}]{
                \texttt{\tiny{[GM$|$GM]}}
                \texttt{```} \\
\texttt{ŘEKNI: Moje pozice: W(3,4), I(3,2), T(6,7), C(2,3), H(2,4). Začnu přesouvat W na (2,2).} \\
\texttt{```} \\
\\ 
\texttt{*(Počkáme na potvrzení druhého hráče a pak provedeme přesun W.)*} \\
            }
        }
    }
    & & \\ \\

    \theutterance \stepcounter{utterance}  
    & & & \multicolumn{2}{p{0.3\linewidth}}{
        \cellcolor[rgb]{0.9,0.9,0.9}{
            \makecell[{{p{\linewidth}}}]{
                \texttt{\tiny{[GM$|$GM]}}
                \texttt{Invalid format: tail is not empty: '} \\
\\ 
\\ 
\texttt{*(Počkáme na potvrzení druhého hráče a pak provedeme přesun W.)*} \\
            }
        }
    }
    & & \\ \\

    \theutterance \stepcounter{utterance}  
    & & \multicolumn{4}{p{0.6\linewidth}}{
        \cellcolor[rgb]{0.9,0.9,0.9}{
            \makecell[{{p{\linewidth}}}]{
                \texttt{\tiny{[P1$\langle$GM]}}
                \texttt{Trest: Vaše zpráva nesmí obsahovat nic po příkazu!} \\
\\ 
\texttt{Ujistěte se, že vaše odpověď obsahuje pouze `ŘEKNI: <ZPRÁVA>` nebo `POHYB: <OBJEKT>, (<X>, <Y>)` a nic jiného!} \\
\\ 
\texttt{Dohromady jste nasbírali 4 z 12 pokut.} \\
\texttt{Zkuste to prosím znovu!} \\
            }
        }
    }
    & \\ \\

    \theutterance \stepcounter{utterance}  
    & \multicolumn{4}{p{0.6\linewidth}}{
        \cellcolor[rgb]{0.9,0.9,0.9}{
            \makecell[{{p{\linewidth}}}]{
                \texttt{\tiny{[P1$\rangle$GM]}}
                \texttt{```} \\
\texttt{ŘEKNI: Moje pozice: W(3,4), I(3,2), T(6,7), C(2,3), H(2,4)} \\
\texttt{```} \\
            }
        }
    }
    & & \\ \\

    \theutterance \stepcounter{utterance}  
    & & & \multicolumn{2}{p{0.3\linewidth}}{
        \cellcolor[rgb]{0.9,0.9,0.9}{
            \makecell[{{p{\linewidth}}}]{
                \texttt{\tiny{[GM$|$GM]}}
                \texttt{```} \\
\texttt{ŘEKNI: Moje pozice: W(3,4), I(3,2), T(6,7), C(2,3), H(2,4)} \\
\texttt{```} \\
            }
        }
    }
    & & \\ \\

    \theutterance \stepcounter{utterance}  
    & & \multicolumn{4}{p{0.6\linewidth}}{
        \cellcolor[rgb]{0.9,0.9,0.9}{
            \makecell[{{p{\linewidth}}}]{
                \texttt{\tiny{[GM$\rangle$P2]}}
                \texttt{Vaše zpráva byla předána druhému hráči.} \\
\\ 
\\ 
\texttt{Právě hrajete kolo 2 z maximálně 20.} \\
\texttt{Dohromady jste nasbírali 4 z 12 pokut.} \\
\texttt{Druhý hráč poslal tuto zprávu:} \\
\\ 
\\ 
\texttt{"Moje pozice: W(3,4), I(3,2), T(6,7), C(2,3), H(2,4)"} \\
\\ 
\\ 
\texttt{Jaký je váš další příkaz?} \\
            }
        }
    }
    & \\ \\

    \theutterance \stepcounter{utterance}  
    & & & \multicolumn{4}{p{0.6\linewidth}}{
        \cellcolor[rgb]{0.9,0.9,0.9}{
            \makecell[{{p{\linewidth}}}]{
                \texttt{\tiny{[GM$\langle$P2]}}
                \texttt{```} \\
\texttt{POHYB: W, (2,2)} \\
\texttt{```} \\
\\ 
\texttt{*(Přesouvám W na navrženou cílovou pozici (2,2), protože je volná a souhlasí s dohodou. Další kroky budou záviset na odpovědi druhého hráče.)*} \\
            }
        }
    }
     \\ \\

    \theutterance \stepcounter{utterance}  
    & & & \multicolumn{2}{p{0.3\linewidth}}{
        \cellcolor[rgb]{0.9,0.9,0.9}{
            \makecell[{{p{\linewidth}}}]{
                \texttt{\tiny{[GM$|$GM]}}
                \texttt{```} \\
\texttt{POHYB: W, (2,2)} \\
\texttt{```} \\
\\ 
\texttt{*(Přesouvám W na navrženou cílovou pozici (2,2), protože je volná a souhlasí s dohodou. Další kroky budou záviset na odpovědi druhého hráče.)*} \\
            }
        }
    }
    & & \\ \\

    \theutterance \stepcounter{utterance}  
    & & & \multicolumn{2}{p{0.3\linewidth}}{
        \cellcolor[rgb]{0.9,0.9,0.9}{
            \makecell[{{p{\linewidth}}}]{
                \texttt{\tiny{[GM$|$GM]}}
                \texttt{Invalid format: tail is not empty: '} \\
\\ 
\\ 
\texttt{*(Přesouvám W na navrženou cílovou pozici (2,2), protože je volná a souhlasí s dohodou. Další kroky budou záviset na odpovědi druhého hráče.)*} \\
            }
        }
    }
    & & \\ \\

    \theutterance \stepcounter{utterance}  
    & & \multicolumn{4}{p{0.6\linewidth}}{
        \cellcolor[rgb]{0.9,0.9,0.9}{
            \makecell[{{p{\linewidth}}}]{
                \texttt{\tiny{[GM$\rangle$P2]}}
                \texttt{Trest: Vaše zpráva nesmí obsahovat nic po příkazu!} \\
\\ 
\texttt{Ujistěte se, že vaše odpověď obsahuje pouze `ŘEKNI: <ZPRÁVA>` nebo `POHYB: <OBJEKT>, (<X>, <Y>)` a nic jiného!} \\
\\ 
\texttt{Dohromady jste nasbírali 5 z 12 pokut.} \\
\texttt{Zkuste to prosím znovu!} \\
            }
        }
    }
    & \\ \\

    \theutterance \stepcounter{utterance}  
    & & & \multicolumn{4}{p{0.6\linewidth}}{
        \cellcolor[rgb]{0.9,0.9,0.9}{
            \makecell[{{p{\linewidth}}}]{
                \texttt{\tiny{[GM$\langle$P2]}}
                \texttt{```} \\
\texttt{POHYB: W, (2,2)} \\
\texttt{```} \\
            }
        }
    }
     \\ \\

    \theutterance \stepcounter{utterance}  
    & & & \multicolumn{2}{p{0.3\linewidth}}{
        \cellcolor[rgb]{0.9,0.9,0.9}{
            \makecell[{{p{\linewidth}}}]{
                \texttt{\tiny{[GM$|$GM]}}
                \texttt{```} \\
\texttt{POHYB: W, (2,2)} \\
\texttt{```} \\
            }
        }
    }
    & & \\ \\

    \theutterance \stepcounter{utterance}  
    & & & \multicolumn{2}{p{0.3\linewidth}}{
        \cellcolor[rgb]{0.9,0.9,0.9}{
            \makecell[{{p{\linewidth}}}]{
                \texttt{\tiny{[GM$|$GM]}}
                \texttt{Penalizace: (2,2) není prázdný, ale obsahuje '│'.} \\
\\ 
\texttt{Dohromady jste nasbírali 6 z 12 pokut.} \\
\texttt{Zkuste to prosím znovu!} \\
            }
        }
    }
    & & \\ \\

    \theutterance \stepcounter{utterance}  
    & & \multicolumn{4}{p{0.6\linewidth}}{
        \cellcolor[rgb]{0.9,0.9,0.9}{
            \makecell[{{p{\linewidth}}}]{
                \texttt{\tiny{[GM$\rangle$P2]}}
                \texttt{Penalizace: (2,2) není prázdný, ale obsahuje '│'.} \\
\\ 
\texttt{Dohromady jste nasbírali 6 z 12 pokut.} \\
\texttt{Zkuste to prosím znovu!} \\
            }
        }
    }
    & \\ \\

    \theutterance \stepcounter{utterance}  
    & & & \multicolumn{4}{p{0.6\linewidth}}{
        \cellcolor[rgb]{0.9,0.9,0.9}{
            \makecell[{{p{\linewidth}}}]{
                \texttt{\tiny{[GM$\langle$P2]}}
                \texttt{```} \\
\texttt{POHYB: W, (2,1)} \\
\texttt{```} \\
\\ 
\texttt{*(Pokusím se přesunout W na sousední volné pole (2,1), abych se vyhnul penalizaci. Pozice (2,2) je blokovaná svislou čarou "│".)*} \\
            }
        }
    }
     \\ \\

    \theutterance \stepcounter{utterance}  
    & & & \multicolumn{2}{p{0.3\linewidth}}{
        \cellcolor[rgb]{0.9,0.9,0.9}{
            \makecell[{{p{\linewidth}}}]{
                \texttt{\tiny{[GM$|$GM]}}
                \texttt{```} \\
\texttt{POHYB: W, (2,1)} \\
\texttt{```} \\
\\ 
\texttt{*(Pokusím se přesunout W na sousední volné pole (2,1), abych se vyhnul penalizaci. Pozice (2,2) je blokovaná svislou čarou "│".)*} \\
            }
        }
    }
    & & \\ \\

    \theutterance \stepcounter{utterance}  
    & & & \multicolumn{2}{p{0.3\linewidth}}{
        \cellcolor[rgb]{0.9,0.9,0.9}{
            \makecell[{{p{\linewidth}}}]{
                \texttt{\tiny{[GM$|$GM]}}
                \texttt{Invalid format: tail is not empty: '} \\
\\ 
\\ 
\texttt{*(Pokusím se přesunout W na sousední volné pole (2,1), abych se vyhnul penalizaci. Pozice (2,2) je blokovaná svislou čarou "│".)*} \\
            }
        }
    }
    & & \\ \\

    \theutterance \stepcounter{utterance}  
    & & \multicolumn{4}{p{0.6\linewidth}}{
        \cellcolor[rgb]{0.9,0.9,0.9}{
            \makecell[{{p{\linewidth}}}]{
                \texttt{\tiny{[GM$\rangle$P2]}}
                \texttt{Trest: Vaše zpráva nesmí obsahovat nic po příkazu!} \\
\\ 
\texttt{Ujistěte se, že vaše odpověď obsahuje pouze `ŘEKNI: <ZPRÁVA>` nebo `POHYB: <OBJEKT>, (<X>, <Y>)` a nic jiného!} \\
\\ 
\texttt{Dohromady jste nasbírali 7 z 12 pokut.} \\
\texttt{Zkuste to prosím znovu!} \\
            }
        }
    }
    & \\ \\

    \theutterance \stepcounter{utterance}  
    & & & \multicolumn{4}{p{0.6\linewidth}}{
        \cellcolor[rgb]{0.9,0.9,0.9}{
            \makecell[{{p{\linewidth}}}]{
                \texttt{\tiny{[GM$\langle$P2]}}
                \texttt{```} \\
\texttt{POHYB: W, (2,1)} \\
\texttt{```} \\
            }
        }
    }
     \\ \\

    \theutterance \stepcounter{utterance}  
    & & & \multicolumn{2}{p{0.3\linewidth}}{
        \cellcolor[rgb]{0.9,0.9,0.9}{
            \makecell[{{p{\linewidth}}}]{
                \texttt{\tiny{[GM$|$GM]}}
                \texttt{```} \\
\texttt{POHYB: W, (2,1)} \\
\texttt{```} \\
            }
        }
    }
    & & \\ \\

    \theutterance \stepcounter{utterance}  
    & & & \multicolumn{2}{p{0.3\linewidth}}{
        \cellcolor[rgb]{0.9,0.9,0.9}{
            \makecell[{{p{\linewidth}}}]{
                \texttt{\tiny{[GM$|$GM]}}
                \texttt{Penalizace: (2,1) není prázdný, ale obsahuje '│'.} \\
\\ 
\texttt{Dohromady jste nasbírali 8 z 12 pokut.} \\
\texttt{Zkuste to prosím znovu!} \\
            }
        }
    }
    & & \\ \\

    \theutterance \stepcounter{utterance}  
    & & \multicolumn{4}{p{0.6\linewidth}}{
        \cellcolor[rgb]{0.9,0.9,0.9}{
            \makecell[{{p{\linewidth}}}]{
                \texttt{\tiny{[GM$\rangle$P2]}}
                \texttt{Penalizace: (2,1) není prázdný, ale obsahuje '│'.} \\
\\ 
\texttt{Dohromady jste nasbírali 8 z 12 pokut.} \\
\texttt{Zkuste to prosím znovu!} \\
            }
        }
    }
    & \\ \\

    \theutterance \stepcounter{utterance}  
    & & & \multicolumn{4}{p{0.6\linewidth}}{
        \cellcolor[rgb]{0.9,0.9,0.9}{
            \makecell[{{p{\linewidth}}}]{
                \texttt{\tiny{[GM$\langle$P2]}}
                \texttt{```} \\
\texttt{ŘEKNI: W nemůžu přesunout na (2,2) ani (2,1) kvůli čáře. Navrhuji W na (3,1) {-} je to volné?} \\
\texttt{```} \\
\\ 
\texttt{*(Zjišťuji alternativní volnou pozici pro W, abychom mohli pokračovat v koordinaci.)*} \\
            }
        }
    }
     \\ \\

    \theutterance \stepcounter{utterance}  
    & & & \multicolumn{2}{p{0.3\linewidth}}{
        \cellcolor[rgb]{0.9,0.9,0.9}{
            \makecell[{{p{\linewidth}}}]{
                \texttt{\tiny{[GM$|$GM]}}
                \texttt{```} \\
\texttt{ŘEKNI: W nemůžu přesunout na (2,2) ani (2,1) kvůli čáře. Navrhuji W na (3,1) {-} je to volné?} \\
\texttt{```} \\
\\ 
\texttt{*(Zjišťuji alternativní volnou pozici pro W, abychom mohli pokračovat v koordinaci.)*} \\
            }
        }
    }
    & & \\ \\

    \theutterance \stepcounter{utterance}  
    & & & \multicolumn{2}{p{0.3\linewidth}}{
        \cellcolor[rgb]{0.9,0.9,0.9}{
            \makecell[{{p{\linewidth}}}]{
                \texttt{\tiny{[GM$|$GM]}}
                \texttt{Invalid format: tail is not empty: '} \\
\\ 
\\ 
\texttt{*(Zjišťuji alternativní volnou pozici pro W, abychom mohli pokračovat v koordinaci.)*} \\
            }
        }
    }
    & & \\ \\

    \theutterance \stepcounter{utterance}  
    & & \multicolumn{4}{p{0.6\linewidth}}{
        \cellcolor[rgb]{0.9,0.9,0.9}{
            \makecell[{{p{\linewidth}}}]{
                \texttt{\tiny{[GM$\rangle$P2]}}
                \texttt{Trest: Vaše zpráva nesmí obsahovat nic po příkazu!} \\
\\ 
\texttt{Ujistěte se, že vaše odpověď obsahuje pouze `ŘEKNI: <ZPRÁVA>` nebo `POHYB: <OBJEKT>, (<X>, <Y>)` a nic jiného!} \\
\\ 
\texttt{Dohromady jste nasbírali 9 z 12 pokut.} \\
\texttt{Zkuste to prosím znovu!} \\
            }
        }
    }
    & \\ \\

    \theutterance \stepcounter{utterance}  
    & & & \multicolumn{4}{p{0.6\linewidth}}{
        \cellcolor[rgb]{0.9,0.9,0.9}{
            \makecell[{{p{\linewidth}}}]{
                \texttt{\tiny{[GM$\langle$P2]}}
                \texttt{```} \\
\texttt{ŘEKNI: Je (3,1) volné pro W?} \\
\texttt{```} \\
            }
        }
    }
     \\ \\

    \theutterance \stepcounter{utterance}  
    & & & \multicolumn{2}{p{0.3\linewidth}}{
        \cellcolor[rgb]{0.9,0.9,0.9}{
            \makecell[{{p{\linewidth}}}]{
                \texttt{\tiny{[GM$|$GM]}}
                \texttt{```} \\
\texttt{ŘEKNI: Je (3,1) volné pro W?} \\
\texttt{```} \\
            }
        }
    }
    & & \\ \\

    \theutterance \stepcounter{utterance}  
    & & \multicolumn{4}{p{0.6\linewidth}}{
        \cellcolor[rgb]{0.9,0.9,0.9}{
            \makecell[{{p{\linewidth}}}]{
                \texttt{\tiny{[P1$\langle$GM]}}
                \texttt{Vaše zpráva byla předána druhému hráči.} \\
\\ 
\\ 
\texttt{Právě hrajete kolo 2 z maximálně 20.} \\
\texttt{Dohromady jste nasbírali 9 z 12 pokut.} \\
\texttt{Druhý hráč poslal tuto zprávu:} \\
\\ 
\\ 
\texttt{"Je (3,1) volné pro W?"} \\
\\ 
\\ 
\texttt{Jaký je váš další příkaz?} \\
            }
        }
    }
    & \\ \\

    \theutterance \stepcounter{utterance}  
    & \multicolumn{4}{p{0.6\linewidth}}{
        \cellcolor[rgb]{0.9,0.9,0.9}{
            \makecell[{{p{\linewidth}}}]{
                \texttt{\tiny{[P1$\rangle$GM]}}
                \texttt{```} \\
\texttt{POHYB: W, (3,1)} \\
\texttt{```} \\
\\ 
\texttt{*(Ověřím, zda je (3,1) volné – pokud ne, dostanu penalizaci, ale druhý hráč se ptal, tak předpokládám, že ano. Pokud ne, upravím další tah.)*} \\
            }
        }
    }
    & & \\ \\

    \theutterance \stepcounter{utterance}  
    & & & \multicolumn{2}{p{0.3\linewidth}}{
        \cellcolor[rgb]{0.9,0.9,0.9}{
            \makecell[{{p{\linewidth}}}]{
                \texttt{\tiny{[GM$|$GM]}}
                \texttt{```} \\
\texttt{POHYB: W, (3,1)} \\
\texttt{```} \\
\\ 
\texttt{*(Ověřím, zda je (3,1) volné – pokud ne, dostanu penalizaci, ale druhý hráč se ptal, tak předpokládám, že ano. Pokud ne, upravím další tah.)*} \\
            }
        }
    }
    & & \\ \\

    \theutterance \stepcounter{utterance}  
    & & & \multicolumn{2}{p{0.3\linewidth}}{
        \cellcolor[rgb]{0.9,0.9,0.9}{
            \makecell[{{p{\linewidth}}}]{
                \texttt{\tiny{[GM$|$GM]}}
                \texttt{Invalid format: tail is not empty: '} \\
\\ 
\\ 
\texttt{*(Ověřím, zda je (3,1) volné – pokud ne, dostanu penalizaci, ale druhý hráč se ptal, tak předpokládám, že ano. Pokud ne, upravím další tah.)*} \\
            }
        }
    }
    & & \\ \\

    \theutterance \stepcounter{utterance}  
    & & \multicolumn{4}{p{0.6\linewidth}}{
        \cellcolor[rgb]{0.9,0.9,0.9}{
            \makecell[{{p{\linewidth}}}]{
                \texttt{\tiny{[P1$\langle$GM]}}
                \texttt{Trest: Vaše zpráva nesmí obsahovat nic po příkazu!} \\
\\ 
\texttt{Ujistěte se, že vaše odpověď obsahuje pouze `ŘEKNI: <ZPRÁVA>` nebo `POHYB: <OBJEKT>, (<X>, <Y>)` a nic jiného!} \\
\\ 
\texttt{Dohromady jste nasbírali 10 z 12 pokut.} \\
\texttt{Zkuste to prosím znovu!} \\
            }
        }
    }
    & \\ \\

    \theutterance \stepcounter{utterance}  
    & \multicolumn{4}{p{0.6\linewidth}}{
        \cellcolor[rgb]{0.9,0.9,0.9}{
            \makecell[{{p{\linewidth}}}]{
                \texttt{\tiny{[P1$\rangle$GM]}}
                \texttt{```} \\
\texttt{POHYB: W, (3,1)} \\
\texttt{```} \\
            }
        }
    }
    & & \\ \\

    \theutterance \stepcounter{utterance}  
    & & & \multicolumn{2}{p{0.3\linewidth}}{
        \cellcolor[rgb]{0.9,0.9,0.9}{
            \makecell[{{p{\linewidth}}}]{
                \texttt{\tiny{[GM$|$GM]}}
                \texttt{```} \\
\texttt{POHYB: W, (3,1)} \\
\texttt{```} \\
            }
        }
    }
    & & \\ \\

    \theutterance \stepcounter{utterance}  
    & & & \multicolumn{2}{p{0.3\linewidth}}{
        \cellcolor[rgb]{0.9,0.9,0.9}{
            \makecell[{{p{\linewidth}}}]{
                \texttt{\tiny{[GM$|$GM]}}
                \texttt{W moved to (3, 1)} \\
            }
        }
    }
    & & \\ \\

    \theutterance \stepcounter{utterance}  
    & & \multicolumn{4}{p{0.6\linewidth}}{
        \cellcolor[rgb]{0.9,0.9,0.9}{
            \makecell[{{p{\linewidth}}}]{
                \texttt{\tiny{[GM$\rangle$P2]}}
                \texttt{Vaše zpráva byla předána druhému hráči.} \\
\\ 
\\ 
\texttt{Právě hrajete kolo 3 z maximálně 20.} \\
\texttt{Dohromady jste nasbírali 10 z 12 pokut.} \\
\texttt{Druhý hráč přesunul objekt na své mřížce.} \\
\\ 
\\ 
\texttt{Jaký je váš další příkaz?} \\
            }
        }
    }
    & \\ \\

    \theutterance \stepcounter{utterance}  
    & & & \multicolumn{4}{p{0.6\linewidth}}{
        \cellcolor[rgb]{0.9,0.9,0.9}{
            \makecell[{{p{\linewidth}}}]{
                \texttt{\tiny{[GM$\langle$P2]}}
                \texttt{```} \\
\texttt{POHYB: I, (3,2)} \\
\texttt{```} \\
\\ 
\texttt{*(Přesouvám I na cílovou pozici (3,2), která je podle dohody volná a druhý hráč tam má svůj I také. Tím se naše pozice srovnají pro tento objekt.)*} \\
            }
        }
    }
     \\ \\

    \theutterance \stepcounter{utterance}  
    & & & \multicolumn{2}{p{0.3\linewidth}}{
        \cellcolor[rgb]{0.9,0.9,0.9}{
            \makecell[{{p{\linewidth}}}]{
                \texttt{\tiny{[GM$|$GM]}}
                \texttt{```} \\
\texttt{POHYB: I, (3,2)} \\
\texttt{```} \\
\\ 
\texttt{*(Přesouvám I na cílovou pozici (3,2), která je podle dohody volná a druhý hráč tam má svůj I také. Tím se naše pozice srovnají pro tento objekt.)*} \\
            }
        }
    }
    & & \\ \\

    \theutterance \stepcounter{utterance}  
    & & & \multicolumn{2}{p{0.3\linewidth}}{
        \cellcolor[rgb]{0.9,0.9,0.9}{
            \makecell[{{p{\linewidth}}}]{
                \texttt{\tiny{[GM$|$GM]}}
                \texttt{Invalid format: tail is not empty: '} \\
\\ 
\\ 
\texttt{*(Přesouvám I na cílovou pozici (3,2), která je podle dohody volná a druhý hráč tam má svůj I také. Tím se naše pozice srovnají pro tento objekt.)*} \\
            }
        }
    }
    & & \\ \\

    \theutterance \stepcounter{utterance}  
    & & \multicolumn{4}{p{0.6\linewidth}}{
        \cellcolor[rgb]{0.9,0.9,0.9}{
            \makecell[{{p{\linewidth}}}]{
                \texttt{\tiny{[GM$\rangle$P2]}}
                \texttt{Trest: Vaše zpráva nesmí obsahovat nic po příkazu!} \\
\\ 
\texttt{Ujistěte se, že vaše odpověď obsahuje pouze `ŘEKNI: <ZPRÁVA>` nebo `POHYB: <OBJEKT>, (<X>, <Y>)` a nic jiného!} \\
\\ 
\texttt{Dohromady jste nasbírali 11 z 12 pokut.} \\
\texttt{Zkuste to prosím znovu!} \\
            }
        }
    }
    & \\ \\

    \theutterance \stepcounter{utterance}  
    & & & \multicolumn{4}{p{0.6\linewidth}}{
        \cellcolor[rgb]{0.9,0.9,0.9}{
            \makecell[{{p{\linewidth}}}]{
                \texttt{\tiny{[GM$\langle$P2]}}
                \texttt{```} \\
\texttt{POHYB: I, (3,2)} \\
\texttt{```} \\
            }
        }
    }
     \\ \\

    \theutterance \stepcounter{utterance}  
    & & & \multicolumn{2}{p{0.3\linewidth}}{
        \cellcolor[rgb]{0.9,0.9,0.9}{
            \makecell[{{p{\linewidth}}}]{
                \texttt{\tiny{[GM$|$GM]}}
                \texttt{```} \\
\texttt{POHYB: I, (3,2)} \\
\texttt{```} \\
            }
        }
    }
    & & \\ \\

    \theutterance \stepcounter{utterance}  
    & & & \multicolumn{2}{p{0.3\linewidth}}{
        \cellcolor[rgb]{0.9,0.9,0.9}{
            \makecell[{{p{\linewidth}}}]{
                \texttt{\tiny{[GM$|$GM]}}
                \texttt{I moved to (3, 2)} \\
            }
        }
    }
    & & \\ \\

    \theutterance \stepcounter{utterance}  
    & & \multicolumn{4}{p{0.6\linewidth}}{
        \cellcolor[rgb]{0.9,0.9,0.9}{
            \makecell[{{p{\linewidth}}}]{
                \texttt{\tiny{[P1$\langle$GM]}}
                \texttt{Úspěšně přesunuto 'W' do (3,1). Vaše aktualizovaná mřížka vypadá takto:} \\
\texttt{```} \\
\texttt{ 1234567} \\
\texttt{╔═╤═════╗} \\
\texttt{║◌│W◌◌◌◌║ 1} \\
\texttt{║◌│I◌◌◌◌║ 2} \\
\texttt{║C│◌◌◌◌◌║ 3} \\
\texttt{║H│◌◌◌◌◌║ 4} \\
\texttt{╟─┼───┐◌║ 5} \\
\texttt{║◌│◌◌◌├─╢ 6} \\
\texttt{║◌│◌◌◌│T║ 7} \\
\texttt{╚═╧═══╧═╝} \\
\texttt{```} \\
\texttt{Právě hrajete kolo 3 z maximálně 20.} \\
\texttt{Dohromady jste nasbírali 11 z 12 pokut.} \\
\texttt{Druhý hráč přesunul objekt na své mřížce.} \\
\\ 
\\ 
\texttt{Jaký je váš další příkaz?} \\
            }
        }
    }
    & \\ \\

    \theutterance \stepcounter{utterance}  
    & \multicolumn{4}{p{0.6\linewidth}}{
        \cellcolor[rgb]{0.9,0.9,0.9}{
            \makecell[{{p{\linewidth}}}]{
                \texttt{\tiny{[P1$\rangle$GM]}}
                \texttt{```} \\
\texttt{ŘEKNI: W je na (3,1). Mám přesunout I na (3,2) (už tam je), pak C na (5,5). Můžeš přesunout T na (4,4)?} \\
\texttt{```} \\
            }
        }
    }
    & & \\ \\

    \theutterance \stepcounter{utterance}  
    & & & \multicolumn{2}{p{0.3\linewidth}}{
        \cellcolor[rgb]{0.9,0.9,0.9}{
            \makecell[{{p{\linewidth}}}]{
                \texttt{\tiny{[GM$|$GM]}}
                \texttt{```} \\
\texttt{ŘEKNI: W je na (3,1). Mám přesunout I na (3,2) (už tam je), pak C na (5,5). Můžeš přesunout T na (4,4)?} \\
\texttt{```} \\
            }
        }
    }
    & & \\ \\

    \theutterance \stepcounter{utterance}  
    & & \multicolumn{4}{p{0.6\linewidth}}{
        \cellcolor[rgb]{0.9,0.9,0.9}{
            \makecell[{{p{\linewidth}}}]{
                \texttt{\tiny{[GM$\rangle$P2]}}
                \texttt{Úspěšně přesunuto 'I' do (3,2). Vaše aktualizovaná mřížka vypadá takto:} \\
\texttt{```} \\
\texttt{ 1234567} \\
\texttt{╔═╤═════╗} \\
\texttt{║◌│◌◌◌◌◌║ 1} \\
\texttt{║◌│I◌◌H◌║ 2} \\
\texttt{║◌│◌◌C◌◌║ 3} \\
\texttt{║◌│◌◌T◌◌║ 4} \\
\texttt{╟─┼───┐◌║ 5} \\
\texttt{║◌│◌◌◌├─╢ 6} \\
\texttt{║W│◌◌◌│◌║ 7} \\
\texttt{╚═╧═══╧═╝} \\
\texttt{```} \\
\texttt{Právě hrajete kolo 4 z maximálně 20.} \\
\texttt{Dohromady jste nasbírali 11 z 12 pokut.} \\
\texttt{Druhý hráč poslal tuto zprávu:} \\
\\ 
\\ 
\texttt{"W je na (3,1). Mám přesunout I na (3,2) (už tam je), pak C na (5,5). Můžeš přesunout T na (4,4)?"} \\
\\ 
\\ 
\texttt{Jaký je váš další příkaz?} \\
            }
        }
    }
    & \\ \\

    \theutterance \stepcounter{utterance}  
    & & & \multicolumn{4}{p{0.6\linewidth}}{
        \cellcolor[rgb]{0.9,0.9,0.9}{
            \makecell[{{p{\linewidth}}}]{
                \texttt{\tiny{[GM$\langle$P2]}}
                \texttt{```} \\
\texttt{POHYB: T, (4,4)} \\
\texttt{```} \\
\\ 
\texttt{*(Přesouvám T na cílovou pozici (4,4), která je volná a odpovídá dohodě. Tím se srovná další objekt.)*} \\
            }
        }
    }
     \\ \\

    \theutterance \stepcounter{utterance}  
    & & & \multicolumn{2}{p{0.3\linewidth}}{
        \cellcolor[rgb]{0.9,0.9,0.9}{
            \makecell[{{p{\linewidth}}}]{
                \texttt{\tiny{[GM$|$GM]}}
                \texttt{```} \\
\texttt{POHYB: T, (4,4)} \\
\texttt{```} \\
\\ 
\texttt{*(Přesouvám T na cílovou pozici (4,4), která je volná a odpovídá dohodě. Tím se srovná další objekt.)*} \\
            }
        }
    }
    & & \\ \\

    \theutterance \stepcounter{utterance}  
    & & & \multicolumn{2}{p{0.3\linewidth}}{
        \cellcolor[rgb]{0.9,0.9,0.9}{
            \makecell[{{p{\linewidth}}}]{
                \texttt{\tiny{[GM$|$GM]}}
                \texttt{Invalid format: tail is not empty: '} \\
\\ 
\\ 
\texttt{*(Přesouvám T na cílovou pozici (4,4), která je volná a odpovídá dohodě. Tím se srovná další objekt.)*} \\
            }
        }
    }
    & & \\ \\

    \theutterance \stepcounter{utterance}  
    & & \multicolumn{4}{p{0.6\linewidth}}{
        \cellcolor[rgb]{0.9,0.9,0.9}{
            \makecell[{{p{\linewidth}}}]{
                \texttt{\tiny{[GM$\rangle$P2]}}
                \texttt{Trest: Vaše zpráva nesmí obsahovat nic po příkazu!} \\
\\ 
\texttt{Ujistěte se, že vaše odpověď obsahuje pouze `ŘEKNI: <ZPRÁVA>` nebo `POHYB: <OBJEKT>, (<X>, <Y>)` a nic jiného!} \\
\\ 
\texttt{Dohromady jste nasbírali 12 z 12 pokut.} \\
\texttt{Zkuste to prosím znovu!} \\
            }
        }
    }
    & \\ \\

    \theutterance \stepcounter{utterance}  
    & & & \multicolumn{4}{p{0.6\linewidth}}{
        \cellcolor[rgb]{0.9,0.9,0.9}{
            \makecell[{{p{\linewidth}}}]{
                \texttt{\tiny{[GM$\langle$P2]}}
                \texttt{```} \\
\texttt{POHYB: T, (4,4)} \\
\texttt{```} \\
\\ 
\texttt{*(Poslední pokus {-} přesouvám T na (4,4), což je cílová pozice podle dohody. Pokud to nevyjde, prohrajeme kvůli 12 penalizacím.)*} \\
            }
        }
    }
     \\ \\

    \theutterance \stepcounter{utterance}  
    & & & \multicolumn{2}{p{0.3\linewidth}}{
        \cellcolor[rgb]{0.9,0.9,0.9}{
            \makecell[{{p{\linewidth}}}]{
                \texttt{\tiny{[GM$|$GM]}}
                \texttt{```} \\
\texttt{POHYB: T, (4,4)} \\
\texttt{```} \\
\\ 
\texttt{*(Poslední pokus {-} přesouvám T na (4,4), což je cílová pozice podle dohody. Pokud to nevyjde, prohrajeme kvůli 12 penalizacím.)*} \\
            }
        }
    }
    & & \\ \\

    \theutterance \stepcounter{utterance}  
    & & & \multicolumn{2}{p{0.3\linewidth}}{
        \cellcolor[rgb]{0.9,0.9,0.9}{
            \makecell[{{p{\linewidth}}}]{
                \texttt{\tiny{[GM$|$GM]}}
                \texttt{Invalid format: tail is not empty: '} \\
\\ 
\\ 
\texttt{*(Poslední pokus {-} přesouvám T na (4,4), což je cílová pozice podle dohody. Pokud to nevyjde, prohrajeme kvůli 12 penalizacím.)*} \\
            }
        }
    }
    & & \\ \\

    \theutterance \stepcounter{utterance}  
    & & & \multicolumn{2}{p{0.3\linewidth}}{
        \cellcolor[rgb]{0.9,0.9,0.9}{
            \makecell[{{p{\linewidth}}}]{
                \texttt{\tiny{[GM$|$GM]}}
                \texttt{Maximum number of penalties exceeded} \\
            }
        }
    }
    & & \\ \\

    \theutterance \stepcounter{utterance}  
    & & & \multicolumn{2}{p{0.3\linewidth}}{
        \cellcolor[rgb]{0.9,0.9,0.9}{
            \makecell[{{p{\linewidth}}}]{
                \texttt{\tiny{[GM$|$GM]}}
                \texttt{Initial states:} \\
\texttt{```} \\
\texttt{Player 1:} \\
\texttt{ 1234567} \\
\texttt{╔═╤═════╗} \\
\texttt{║◌│◌◌◌◌◌║ 1} \\
\texttt{║◌│I◌◌◌◌║ 2} \\
\texttt{║C│◌◌◌◌◌║ 3} \\
\texttt{║H│◌W◌◌◌║ 4} \\
\texttt{╟─┼───┐◌║ 5} \\
\texttt{║◌│◌◌◌├─╢ 6} \\
\texttt{║◌│◌◌◌│T║ 7} \\
\texttt{╚═╧═══╧═╝} \\
\\ 
\\ 
\texttt{Player 2:} \\
\texttt{ 1234567} \\
\texttt{╔═╤═════╗} \\
\texttt{║◌│◌◌◌◌◌║ 1} \\
\texttt{║◌│◌◌◌H◌║ 2} \\
\texttt{║I│◌◌C◌◌║ 3} \\
\texttt{║◌│◌◌T◌◌║ 4} \\
\texttt{╟─┼───┐◌║ 5} \\
\texttt{║◌│◌◌◌├─╢ 6} \\
\texttt{║W│◌◌◌│◌║ 7} \\
\texttt{╚═╧═══╧═╝} \\
\\ 
\texttt{```} \\
            }
        }
    }
    & & \\ \\

    \theutterance \stepcounter{utterance}  
    & & & \multicolumn{2}{p{0.3\linewidth}}{
        \cellcolor[rgb]{0.9,0.9,0.9}{
            \makecell[{{p{\linewidth}}}]{
                \texttt{\tiny{[GM$|$GM]}}
                \texttt{End states:} \\
\texttt{```} \\
\texttt{Player 1:} \\
\texttt{ 1234567} \\
\texttt{╔═╤═════╗} \\
\texttt{║◌│W◌◌◌◌║ 1} \\
\texttt{║◌│I◌◌◌◌║ 2} \\
\texttt{║C│◌◌◌◌◌║ 3} \\
\texttt{║H│◌◌◌◌◌║ 4} \\
\texttt{╟─┼───┐◌║ 5} \\
\texttt{║◌│◌◌◌├─╢ 6} \\
\texttt{║◌│◌◌◌│T║ 7} \\
\texttt{╚═╧═══╧═╝} \\
\\ 
\\ 
\texttt{Player 2:} \\
\texttt{ 1234567} \\
\texttt{╔═╤═════╗} \\
\texttt{║◌│◌◌◌◌◌║ 1} \\
\texttt{║◌│I◌◌H◌║ 2} \\
\texttt{║◌│◌◌C◌◌║ 3} \\
\texttt{║◌│◌◌T◌◌║ 4} \\
\texttt{╟─┼───┐◌║ 5} \\
\texttt{║◌│◌◌◌├─╢ 6} \\
\texttt{║W│◌◌◌│◌║ 7} \\
\texttt{╚═╧═══╧═╝} \\
\\ 
\texttt{```} \\
            }
        }
    }
    & & \\ \\

    \theutterance \stepcounter{utterance}  
    & & & \multicolumn{2}{p{0.3\linewidth}}{
        \cellcolor[rgb]{0.9,0.9,0.9}{
            \makecell[{{p{\linewidth}}}]{
                \texttt{\tiny{[GM$|$GM]}}
                \texttt{* success: False} \\
\texttt{* lose: False} \\
\texttt{* aborted: True} \\
\texttt{{-}{-}{-}{-}{-}{-}{-}} \\
\texttt{* Shifts: 1.00} \\
\texttt{* Max Shifts: 8.00} \\
\texttt{* Min Shifts: 4.00} \\
\texttt{* End Distance Sum: 19.32} \\
\texttt{* Init Distance Sum: 19.47} \\
\texttt{* Expected Distance Sum: 20.95} \\
\texttt{* Penalties: 13.00} \\
\texttt{* Max Penalties: 12.00} \\
\texttt{* Rounds: 3.00} \\
\texttt{* Max Rounds: 20.00} \\
\texttt{* Object Count: 5.00} \\
            }
        }
    }
    & & \\ \\

\end{supertabular}
}

\end{document}
